\section{The v4.6.8 Measurement Bridge Contract}
\label{sec:measurement468}
% =====================================================================

The v4.6.7 final closure proposal made the theory side complete; what
remains is exactly five artefact-validation groups.  v4.6.8 ships
those five groups as an executable contract
(\texttt{tfpt\_cmb.measurement\_bridge\_v468}) with a deterministic
acceptance gate.

\subsection{The five validation groups}

\begin{center}
\renewcommand{\arraystretch}{1.18}
\begin{tabularx}{\textwidth}{@{}p{0.06\textwidth}p{0.18\textwidth}p{0.27\textwidth}Y@{}}
\toprule
\textbf{ID} & \textbf{Group} & \textbf{Acceptance criterion} & \textbf{Closes (artefact)} \\
\midrule
G1 & \texttt{sprint\_2a\_branch\_bundle} (12 files) &
\(\mathrm{branch\_origin}=\mathrm{real}\) AND 24/24 hard checks of
\texttt{sprint2a\_contract.validate\_sprint2a\_branch} pass.
& E1, E2, E3, E4, E5 \\
G2 & \texttt{sprint\_2c\_thermal\_matter} (7 files) &
QKE Ward residues
\(\mathrm{Tr}_{\mathrm{adm}}(K_{\mathrm{weak}})=c_3\),
\(\mathrm{Tr}_{\mathrm{adm}}(K_{\mathrm{EM}})=\beta\) numerically
verified to \(10^{-5}\); BBN \(Y_p=1/4-\beta\) to \(10^{-4}\); 48
neutral admissible modes enumerated.
& C2, C3, D2 \\
G3 & \texttt{sprint\_2d\_reionization} (5 files) &
\(\tau_{\mathrm{numeric}}=\phi_0+c_3^2/2\) to \(10^{-3}\) relative.
& C4 \\
G4 & \texttt{planck\_likelihood} (4 files) &
Official \texttt{clik} or \texttt{camb\_python\_likelihood};
\(|\Delta(-2\ln\mathcal L)|<5\) vs Planck-2018 \(\Lambda\)CDM;
lensed; no parameter refit.
& --- \\
G5 & \texttt{paper5\_metrology\_corollary} (1 file) &
\texttt{paper5\_metrology\_corollary.tex} contains four formal
identifications: \(\omega_\gamma=\pi^2 c_3^4\),
\(\omega_b=\lambda_C/10\), \(\mathrm{Tr}_{\Lambda^2 E}(P_B\rho_{\mathrm{flavor}})=\Omega_b h^2\),
\(E_\nu^\Sigma=1\,\mathrm{eV}\).
& C1, D1, D3 \\
\bottomrule
\end{tabularx}
\end{center}

In total: 29 required files, 8 numerical checks, 1 acceptance gate.

\subsection{The 12 G1 files}

\begin{center}
\small
\begin{tabular}{ll}
\toprule
File & Purpose \\
\midrule
\texttt{branch\_metadata.json}        & metadata, must declare \(\mathrm{branch\_origin}=\mathrm{real}\) \\
\texttt{B\_sigma\_adm.npz}            & boundary Dirac/seam operator \\
\texttt{P\_adm.npz}                   & admissible projector \(P_{\mathrm{prim}}P_{\mathrm{sing}}P_\Theta\) \\
\texttt{P\_C.npz}                     & Calderon polarisation projector \\
\texttt{iota\_C\_adm.npz}             & Calderon sheet involution \\
\texttt{Delta\_prim\_adm.npz}         & primitive Laplacian \\
\texttt{Q\_top\_adm.npz}              & topological phase generator \\
\texttt{K\_zeta\_adm.npz}             & primordial-curvature Hessian \(K_\Sigma\) \\
\texttt{u\_sigma\_adm.npz}            & primitive seam state, \([u_\Sigma]=1\) \\
\texttt{spectral\_flag\_adm.json}     & canonical flag (no Haar seed) \\
\texttt{observer\_embedding\_planck.json} & tetrad SO(3) Galactic, \(\eta_{\mathrm{LS}}\), parity \\
\texttt{shell\_power\_validation.json} & per-shell \((P_{\mathrm{target}},P_{\mathrm{realized}})\) \\
\bottomrule
\end{tabular}
\end{center}

\subsection{Forbidden inputs}

Every group refuses any artefact whose filename matches
\texttt{planck\_phase}, \texttt{planck\_alm}, \texttt{*\_phase}
(SMICA/NILC/SEVEM/Commander), \texttt{haar\_random\_seed}, or
\texttt{synthetic\_v46}.  The contract is \emph{strict}: even a single
forbidden filename in a group's sub-tree fails that group, regardless
of the numerical-check results.

\subsection{Skeleton emitter and CLI}

\begin{lstlisting}[language=bash]
# Emit a documented skeleton tree (5 sub-directories with READMEs):
python3 scripts/validate_measurement_bridge.py \
  --emit-skeleton input/measurement_bridge_v5

# Validate an existing tree:
python3 scripts/validate_measurement_bridge.py \
  --root input/measurement_bridge_v5

# Smoke test: emit and validate a temp skeleton (FAILS by design):
python3 scripts/validate_measurement_bridge.py --skeleton-and-validate
\end{lstlisting}

The validator emits \texttt{MEASUREMENT\_BRIDGE\_VALIDATION.json} into
the root directory with full per-group, per-file, per-check detail.
Exit code is 0 iff all five groups pass.

\subsection{The acceptance gate}

\begin{lstlisting}[language=Python]
from tfpt_cmb.measurement_bridge_v468 import (
    validate_measurement_bridge, can_flip_to_derived_tfpt,
)

report = validate_measurement_bridge("input/measurement_bridge_v5")
ok, missing = can_flip_to_derived_tfpt(report)

# After v4.6.7 (theory complete) plus v4.6.8 (artefacts validated):
# ok == True  =>  status flips from
#                   derived_tfpt_theory_complete_pending_artifacts
#                 to
#                   derived_tfpt
\end{lstlisting}

\subsection{The complete status flow}

\begin{center}
\renewcommand{\arraystretch}{1.18}
\begin{tabularx}{\textwidth}{@{}p{0.36\textwidth}lY@{}}
\toprule
\textbf{Status} & \textbf{Cycle} & \textbf{What flips} \\
\midrule
\textsf{derived\_tfpt\_candidate\_v46} & v4.6 &
Initial candidate; lemmas A--E proposed, none discharged. \\
\textsf{derived\_tfpt\_theory\_complete\_\allowbreak{}pending\_artifacts}
& v4.6.7 &
\texttt{can\_flip\_to\_theory\_complete(REQUIRED\_DISCHARGE\_OUTCOMES\_V467)} returns \texttt{True}. All 14 hooks resolved by 7 bridge theorems. \\
\textsf{derived\_tfpt} & v4.6.8 (final) &
\texttt{can\_flip\_to\_derived\_tfpt(report)} returns \texttt{True}.
All 5 measurement-bridge groups PASS. \\
\bottomrule
\end{tabularx}
\end{center}

\begin{tfptfrontbox}
\textbf{The bilance after v4.6.8:}\\
\emph{Code complete. Audit gate armed. Proof skeletons complete.
Proof discharge formalised. Theory closure complete. Measurement
bridge contract executable. The path from candidate to
derived\_tfpt requires nothing more than dropping 29 real artefacts
into the bridge tree and running the validator.}
\end{tfptfrontbox}

% =====================================================================
\section{Kill conditions and audit table}
\label{sec:kills}
% =====================================================================

\begin{center}
\renewcommand{\arraystretch}{1.05}
\begin{tabularx}{\textwidth}{@{}p{0.26\textwidth}Y@{}}
\toprule
\textbf{Module} & \textbf{Kill condition} \\
\midrule
\texttt{ledger}        & \(\xi\) audit mismatch unresolved when used in production. \\
\texttt{native\_closure\_v46} & Any of the seven primitives violated; any of Lemmas A--D falsified. \\
\texttt{branch\_phase\_selection} & \(K_\Sigma\) not Hermitian; \(\iota_C\) not an involution; any input matching \texttt{planck\_phase}, \texttt{planck\_alm}, \texttt{smica\_alm\_as\_phase\_seed}, \texttt{nilc\_phase}, \texttt{sevem\_phase}, \texttt{commander\_phase} (Mode E refused). \\
\texttt{birefringence} & Round-trip rotation by \(\pm\beta\) not returning the input within machine precision. \\
\texttt{scalaron}      & \(\alpha_R^\Sigma\le 0\); \(A_s\le 0\); \(n_s\notin(0.5,1.2)\); \(r<0\). \\
\texttt{thermal}       & \(T_{\mathrm{reh}}<1\,\mathrm{MeV}\). \\
\texttt{baryogenesis}  & \(\eta_B\le 0\) or \(\eta_B>10^{-7}\). \\
\texttt{neutrino}      & PMNS unitarity residual \(>10^{-12}\); negative individual mass. \\
\texttt{axion}         & \(\Omega_a h^2\) cannot be saturated within the QCD axion window. \\
\texttt{background}    & \(h\notin(0.4,1.0)\); non-flat with \(|\Omega_K|>10^{-3}\). \\
\texttt{boltzmann}     & \texttt{TFPTCosmology.validate()} non-empty. \\
\texttt{maps}          & Wrong HEALPix ordering or polarisation convention. \\
\texttt{compare}       & A claim of map-level pixel agreement absent the conditioned reconstruction layer. \\
\texttt{run\_cmb\_stage} & \texttt{--mode E} called without \texttt{--candidate-v46}; \texttt{derived\_tfpt\_candidate\_v46} accepted without \texttt{--candidate-v46}; strict-map mode mixed with \texttt{--allow-tau-nuisance}. \\
\texttt{sprint2a\_contract} & Any of the ten hard checks failing on a bundle whose \texttt{branch\_origin.json::origin} is set to ``real''; any \texttt{planck\_phase}, \texttt{planck\_alm}, \texttt{*\_phase} or \texttt{haar\_random\_seed} pattern inside the bundle directory. \\
\bottomrule
\end{tabularx}
\end{center}

% =====================================================================
\section{What this pipeline does \emph{not} claim}
\label{sec:nonclaims}
% =====================================================================

\begin{itemize}[leftmargin=1.4em]
\item \textbf{No theorem-level lemmas.}  Lemmas A--E are conjectures
with proof skeletons (\S\ref{sec:lemmas}); none has a finished proof
at TFPT v4.6.  The pipeline runs in \emph{candidate} mode only.
\item \textbf{No real Mode-E branch.}  The Mode-E run uses the
\textbf{candidate v4.6 synthetic} bundle, \emph{not} a real
Sprint-2A extraction of \(K_\Sigma\), \(u_\Sigma\), \(\iota_C\) from
the boundary data.  Real Mode E is a Sprint 2A--2D deliverable.
\item \textbf{No claim of pixel-level Planck agreement.}  Mode~A is a
random TFPT realisation, by construction independent of Planck phases.
Mode~E phases come from \(\arg\langle e_j,u_\Sigma\rangle\), which
is not the Planck realisation \emph{unless} the real
\(K_\Sigma\) is supplied; the synthetic bundle is structural, not
realisation-aware.
\item \textbf{No Likelihood ``win'' against Planck.}  The TT
\(\chi^2/\mathrm{dof}\) is diagonal-Knox, not the official Planck
likelihood with its full covariance, foreground marginalisation, and
beam/noise model.  TE/EE/BB/TB/EB on the component-separated maps are
\emph{diagnostic only} per rules R\(_2\)/R\(_4\).
\item \textbf{No claim that the Eskilt+22 birefringence
\(0.34^\circ\) is the TFPT \(\beta\).}  The TFPT prediction is
\(0.2424^\circ\); the in-pipeline EB on SMICA gives
\(0.027^\circ\); the discrepancy is dominated by foreground/noise
residuals on a single component map.
\item \textbf{No anchoring on COBE / FIRAS / Planck.}  Strict-map mode
explicitly forbids it; the runner refuses upstream JSONs whose
\textsf{status\_flags} contain \texttt{calibration},
\texttt{nuisance}, \texttt{operational\_closure\_v1},
\texttt{audit\_mismatch}, or \texttt{derived\_fixture}.
\end{itemize}

% =====================================================================
\section{Roadmap to a finished v5.0}
\label{sec:roadmap}
% =====================================================================

The five outstanding theorem-level closures are
\begin{enumerate}[leftmargin=1.4em]
\item \textbf{Lemma A} (Seam Determinant): rigorous definition of
\(\det_{\mathrm{adm}}^{\mathrm{prim}}\) and proof that
non-primitive admissible repetitions vanish.
\item \textbf{Lemma B} (Curvature Residue): proof of
\(\chi_R = (4\pi/\Omega_{\mathrm{adm}})\,e^{-1/\phi_0}\) from the
admissible curvature path-integral.
\item \textbf{Lemma C}: subdivide into C1--C4; C4 (\(\tau_{\mathrm{reio}}\))
is the weakest link --- needs SFR/\(N_\gamma\)/\(f_{\mathrm{esc}}\) to
collapse onto \(\phi_0+c_3^2/2\) without astrophysical fitting.
\item \textbf{Lemma D}: proof of the carrier traces plus the
mass-metrology bridge identifying the dimensionless Takagi trace with
eV.
\item \textbf{Lemma E}: existence and uniqueness of a canonical
spectral flag for \(K_\Sigma\); seam orientation convention; observer
embedding; coupling to CAMB transfer functions.
\end{enumerate}

The pipeline is designed so that, once any of these proofs lands, the
status flag flips from \textsf{derived\_tfpt\_candidate\_v46} to
\textsf{derived\_tfpt} without changing any downstream module.  The
Sprint-2A contract (\S\ref{sec:sprint2a}) provides the executable
acceptance gate: the moment a real boundary extraction passes all
24/24 hard checks with \(\mathrm{branch\_origin}=\mathrm{real}\), the
status flag of every Mode-E run flips from
\textsf{lemma\_E\_candidate\_v46} to
\textsf{lemma\_E\_proven}.  No further code changes are required.

\paragraph{Out-of-sample tests.} Even before the proofs land, the
following are testable predictions of v4.6 and would either
falsify the conjecture or sharpen the agreement:
\begin{itemize}[leftmargin=1.4em]
\item BAO scale (DESI Y3, eBOSS) at 0.5\% precision;
\item CMB lensing power spectrum (Planck PR4, ACT DR6, SPT-3G);
\item \(S_8\) tension (KiDS, DES, Euclid);
\item Deuterium primordial abundance (high-resolution QSO absorption);
\item Effective neutrino number from BBN (independent of CMB);
\item Direct \(\sum m_\nu\) (\,KATRIN, Project~8, JUNO\,);
\item Tensor-to-scalar ratio \(r=3.45\times 10^{-3}\) (LiteBIRD,
PICO);
\item 21~cm reionization signal (HERA, SKA-Low);
\item Numerology grammar test: enumerate every reasonable
algebraic combination of the seven primitives; count how many also
land within \(2\sigma\) of Planck for the 26 parameters.  This is the
sharpest single defence against ``hübsche Zahlenspielerei.''
\end{itemize}

\begin{tfptexports}
\textbf{Code-side complete.} Exports: the primitive ledger JSON, the
26-parameter \texttt{cosmology\_candidate.json} (every field
\textsf{derived\_tfpt\_candidate\_v46}), the lensed CMB spectra
\(C_\ell^{XY}\), Mode-A HEALPix maps at \(N_{\mathrm{side}}=1024\),
Mode-E HEALPix maps at \(N_{\mathrm{side}}=512\) with the Lemma-E
phases, the per-component Planck PR3 comparison report, the
five-figure comparison suite, and the new v4.6.4 Sprint-2A contract
validator.  v4.6.5 adds the full Lemma proof outlines
(\texttt{tfpt\_cmb.lemma\_proof\_outlines}) with 14 hook solution
sketches.

\textbf{Data-side incomplete.} Real \(K_\Sigma\), \(u_\Sigma\),
\(\iota_C\) extraction from a Sprint-2A boundary run; QKE collision
kernels for \(N_{\mathrm{eff}}\); BBN reaction amplitudes from
\(\Gamma_{\mathrm{adm,k}}\); reionization Boltzmann from the TFPT
ionising history.

\textbf{Proof-side incomplete.} 14 hooks across Lemmas A--E, all
formally documented in Appendix~\ref{app:hooks} with proposed
solutions.
\end{tfptexports}

% =====================================================================
\section{v4.10--v4.12: Operator-Discharged Closure and the Planck Pass}
\label{sec:v412}
% =====================================================================

The v4.6 cycle delivered the \emph{candidate} closure with five lemmas
A--E and a 14-hook proof discharge.  The next two sprints (v4.7, v4.8)
narrowed the geometry to the Lens space \(L(10,1)\times\mathbb Z_2\)
on the spectral-action corridor, and v4.9's Fisher residual
decomposition pinpointed the remaining \(\Delta(-2\ln L)=+6.04\) gap
to two parameters: \(\log A\) and \(\tau_{\mathrm{reio}}\).  Sprints
v4.10--v4.12 closed this gap operatorically, without touching any
upstream Lemma~A--D module.  This section summarises the result.

\subsection{Sprint history v4.6 \(\to\) v4.12}
\label{sec:v412_history}

Figure~\ref{fig:v412_sprint_history} shows the
\(\Delta(-2\ln L)\) value vs Planck~2018 \(\Lambda\)CDM across all
sprints between v4.6 and v4.12.  The strict TFPT mode never refits
parameters and always uses \texttt{A\_planck = 1.0} on the official
Plik-lite native + lowl TT/EE likelihood (Cobaya).  The drop from
v4.9 (\(+6.04\)) to v4.10/v4.11/v4.12 (\(+0.61\)/\(+0.6119\)/\(+0.6083\))
is \emph{not} a fit improvement --- it is the operator-derived BT2
correction landing on the Fisher residual direction.

\begin{figure}[!htb]
\centering
\includegraphics[width=\textwidth]{fig12_v412_sprint_history.png}
\caption{Sprint history of the strict-mode TFPT \(\Delta(-2\ln L)\)
against Planck~2018 \(\Lambda\)CDM, from v4.6 tree-level
(\(+189.86\)) through the v4.7.2 hyperbolic No-Go (\(+5754\)),
v4.7.6 Spectral-Action UV identification of Lens \(L(10,1)\)
(\(+8.5\)), v4.8 four-corner corridor hit (\(+6.04\)),
v4.9 Fisher diagnosis, the v4.10 closed-form lemma refinements
(\(+0.61\)), v4.11 BT2 operator discharge (\(+0.6119\)), and the
v4.12 \(W_B\) BSE operator discharge (\(+0.6083\), canonical
\texttt{A\_qcd\_curvature\_gauge}).  Dashed: Stefan's strict gate
\(\Delta < 5\).  No parameter is refit; \texttt{A\_planck = 1.0}
throughout.}
\label{fig:v412_sprint_history}
\end{figure}

\subsection{The visibility-amplitude Ward identity (v4.11)}
\label{sec:v412_ward}

Stefan v4.11 brief Section~3 introduced the central new theorem:

\begin{equation}
\boxed{\;
\Delta_R^\Sigma \;=\; 2\,\Delta\tau_{\mathrm{sheet}}
\;+\; O(\delta_{\mathrm{top}})
\;}
\label{eq:ward_identity}
\end{equation}

with the explicit closed forms

\begin{align}
\Delta_R^\Sigma &= \beta\bigl(1-\tfrac{\lambda_C}{2}\bigr)
\;-\; \delta_{\mathrm{top}}, \qquad
\delta_{\mathrm{top}} = \Omega_{\mathrm{adm}}\,c_3^4 ,\\
\Delta\tau_{\mathrm{sheet}} &= c_3^2\bigl(1+\tfrac{\lambda_C}{2}\bigr).
\end{align}

Numerically:

\begin{center}
\renewcommand{\arraystretch}{1.15}
\begin{tabular}{lr}\hline\hline
quantity & value \\\hline
\(\Delta_R^\Sigma\)                & \(+3.6363\times 10^{-3}\) \\
\(2\,\Delta\tau_{\mathrm{sheet}}\) & \(+3.5215\times 10^{-3}\) \\
\(\delta_{\mathrm{top}}\)          & \(+1.2030\times 10^{-4}\) \\
residual / \(\delta_{\mathrm{top}}\) & \(0.954\) \\
ward identity holds                  & TRUE \\\hline
\end{tabular}
\end{center}

The Ward identity proves that the BT2 magnitude correction and the
Lemma~C reionisation correction describe the \emph{same} physical
direction in the visibility amplitude \(A_s\,e^{-2\tau}\):

\begin{equation}
\Delta\log A_s = -\Delta_R^\Sigma, \qquad
e^{-2\Delta\tau_{\mathrm{sheet}}} \approx
e^{-\Delta_R^\Sigma + O(\delta_{\mathrm{top}})}.
\end{equation}

Consequently applying both is double counting and is explicitly
refused by the master pipeline:

\begin{lstlisting}[language=Python]
if apply_delta_R and apply_delta_tau and not orthogonalisation_present:
    raise DoubleCountingError(
        "BT2 and tau corrections are the same visibility-"
        "amplitude direction A_s exp(-2 tau)"
    )
\end{lstlisting}

The canonical choice is therefore the \emph{curvature gauge}: apply
\(\Delta_R^\Sigma\) to \(Z_R\) and keep \(\tau\) at its
leading-order Lemma C value \(\tau_{\mathrm{LO}}=\phi_0+\tfrac12 c_3^2\).
The visibility gauge (\(\Delta\tau_{\mathrm{sheet}}\) only) is a
candidate dual but \emph{not} an additive correction.

